\section{old Background}\label{sec1}

High-throughput sequencing has enabled the characterization of complex microbial communities across biological, clinical, and environmental contexts. Taxonomic profiling commonly relies on marker-gene amplicon sequencing or shotgun metagenomics \cite{jovel_characterization_2016,durazzi_comparison_2021}. Marker-gene approaches target taxonomically informative regions, most commonly the 16S rRNA gene for bacterial and archaeal communities, to identify operational taxonomic units (OTUs). In contrast, shotgun metagenomics sequences DNA without targeting specific genomic regions, enabling broader taxonomic and functional characterization \cite{jovel_characterization_2016,durazzi_comparison_2021,breitwieser_review_2019}. Metatranscriptomic sequencing similarly profiles microbial communities from RNA, providing information on transcriptionally active microorganisms and their expressed functions \cite{breitwieser_review_2019,bashiardes_use_2016}. Despite their methodological differences, these approaches can generate taxonomic abundance profiles describing the microorganisms detected across multiple samples.

Microbial taxa are organized within a hierarchical taxonomy, represented as a branching tree \cite{letunic_interactive_2021}, in which each organism belongs to increasingly broader taxonomic groups, from species and genera to families, orders, classes, phyla, and domains \cite{woese_towards_1990}. Exploring microbial abundance data therefore involves examining abundance patterns across samples while navigating different levels of the taxonomic hierarchy. As the number of taxa and samples increases, representing and exploring these relationships becomes increasingly challenging \cite{wagner_metaviz_2018}.


The interpretation of microbial abundance profiles depends not only on taxonomic abundances but also on information associated with the samples themselves. Biological and clinical metadata have been widely used to investigate microbiome-host associations, notably through large-scale initiatives such as the Human Microbiome Project (HMP) \cite{huttenhower_structure_2012,proctor_integrative_2019}, with associations reported for type 2 diabetes \citep{qin_metagenome-wide_2012} and cancer immunotherapy response \citep{routy_gut_2018,gopalakrishnan_gut_2018}. Conversely, technical metadata can help identify variation associated with sample collection, processing, library preparation, or sequencing rather than biological differences. For example, microbial contaminants can vary between laboratories, reagents, and reagent batches \cite{salter_reagent_2014,laurence_common_2014}, an issue of particular relevance when microbial reads represent only a small fraction of host-dominated sequencing data \cite{ge_comprehensive_2025}. Joint exploration of taxonomic abundances and sample metadata can therefore help investigate biological or clinical phenotypes while also identifying patterns associated with technical variables.




The exploration of these abundance profiles is further complicated by intrinsic properties of microbiome sequencing data, including differences in sequencing depth across samples, sparsity, compositionality, and highly skewed abundance distributions \cite{weiss_normalization_2017,lin_analysis_2020,llorens-rico_benchmarking_2021}. Numerous normalization and transformation strategies have been proposed to account for these properties, and benchmarking studies have shown that the selected preprocessing strategy can influence downstream analyses, including differential abundance testing, taxon--metadata associations, and phenotype prediction \cite{weiss_normalization_2017,chen_gmpr_2018,llorens-rico_benchmarking_2021,nearing_microbiome_2022,wang_comparison_2024}. The relative performance of these approaches depends on both the characteristics of the dataset and the objective of the analysis, with no single strategy being optimal in every setting. Visualizing abundance profiles under different preprocessing strategies can therefore provide complementary views of the same dataset, allowing researchers to directly examine how these choices affect abundance distributions and identify patterns that are sensitive to the selected method.

Beyond the choice of preprocessing methods, exploring microbial abundance profiles requires examining how taxa are distributed across individual samples. Abundance tables and summary statistics provide quantitative information but may obscure differences between samples or patterns associated with specific taxa. Interactive visualization offers a complementary approach by allowing researchers to progressively explore abundance distributions, focus on subsets of taxa or samples, and investigate relationships with associated metadata. Several tools have been developed to support these exploratory tasks (Table~\ref{tab:tool_comparison}). Krona uses an interactive sunburst chart to navigate hierarchical taxonomic profiles \cite{ondov_interactive_2011}, Metacoder represents taxonomic abundances and differences between groups using heat trees \cite{foster_metacoder_2017}, and Pavian uses Sankey diagrams to represent the distribution of reads across taxonomic levels, complemented by heatmaps for multi-sample comparisons \cite{breitwieser_pavian_2020}. Other tools integrate multiple visualization techniques within broader microbiome analysis environments. Shiny-phyloseq provides bar plots, heatmaps, phylogenetic trees, ordination plots, and interactive network representations \cite{mcmurdie_shiny-phyloseq_2015}. QIIME2 View provides browser-based exploration of QIIME 2 outputs through interactive visualizations including taxonomic bar plots at different taxonomic levels, PCoA ordination plots, alpha-diversity raincloud plots, and summary statistics \cite{bolyen_reproducible_2019}. MicrobiomeAnalyst combines statistical and preprocessing methods with interactive stacked bar plots, heatmaps, and other representations for microbiome data exploration \cite{lu_microbiomeanalyst_2023}. Metaviz combines its FacetZoom representation for navigating the taxonomic hierarchy with coordinated heatmaps, stacked bar plots, and other sample-level visualizations, allowing users to interactively select and highlight corresponding data across multiple views \cite{wagner_metaviz_2018}. Snowflake represents microbiome data as bipartite graphs connecting sample nodes to the microbial features detected within them, allowing the identification of shared microorganisms while incorporating taxonomic annotations and sample metadata \cite{peeters_snowflake_2024}. These tools provide complementary approaches to microbiome visualization but differ in how taxonomic structure, individual samples, associated metadata, and preprocessing choices can be explored together.

Here, we present TaxaLoom, a web-based application for the interactive exploration of sample-level taxonomic abundance profiles. TaxaLoom combines taxonomic hierarchy navigation, sample metadata, synchronized visualizations, and the comparison of preprocessing strategies within a common interface. Although developed using shotgun sequencing data, TaxaLoom can be applied to taxonomic abundance profiles provided as raw count data, independently of the sequencing approach used to generate them. TaxaLoom thereby facilitates the exploration of abundance patterns across taxa, samples, and associated metadata.


\section{à garder}

Unlike DNA-based metagenomic sequencing, where the sequenced material reflects the genomic content of the microbial community, metatranscriptomic profiles are additionally influenced by differences in gene expression (Franzosa et al., 2014).This additional source of variation contributes to the wide range of read counts observed across taxa and can produce strongly right-skewed abundance distributions. Together with differences in sequencing depth, compositionality, and sparsity, these characteristics motivate the use of different normalization and transformation strategies before downstream analysis.


As a reminder, normalization methods aim to make values comparable between samples, while transformation changes the mathematical scale, for example through logarithmic or square root transformations. As many preprocessing methods combine both operations under the broader term of ``normalization'', the two terms will sometimes be used interchangeably in the following sections.


